% CS 6120 Final Project Report (LaTeX)
\documentclass[11pt]{article}
\usepackage[utf8]{inputenc}
\usepackage[margin=1in]{geometry}
\usepackage[hidelinks]{hyperref}
\usepackage{parskip}
\usepackage{enumitem}
\title{Conference Matching Platform: \\ \large Hybrid RAG and Semantic Retrieval for Professional Event Networking}
\author{Xian Cao \\ \texttt{https://github.com/xiancao2024/Conference-Matching-Platform}}
\date{April 2025}

\begin{document}
\maketitle

\section*{1 Objectives and Introduction}

This project builds a retrieval-augmented generation (RAG) matching platform for conference attendees and sessions. Using an event attendance dataset, the system normalizes attendance rows into a conference schema (attendees and session resources), indexes those entities, and runs a hybrid matching pipeline mixing lexical and semantic retrieval. The objective is to demonstrate a real-data RAG / PAL-style assistant that can answer matchmaking queries, surface attendees and sessions, and explain matches.

\section*{2 Background and Related Work}

\textbf{Dataset.} The Kaggle Event Attendance Dataset (cankatsrc/event-attendance-dataset) contains event attendance rows with fields such as Event ID, Event Name, Location, Date \\& Time, Attendee Name, Attendee Email, and Attendee Phone. The dataset is suitable because it contains many examples of real-world event attendance which can be normalized into structured entities for retrieval and evaluation.

\textbf{Related work.} RAG systems for QA and assistant workflows (e.g., Lewis et al., 2020; Karpukhin et al., 2020) and PAL (Program-Aided Language models) approaches where an LLM orchestrates retrieval and programmatic filters. This project adapts hybrid retrieval ideas to matchmaking over imported attendance records.

\section*{3 Approach and Implementation}

\subsection*{Overview}
\begin{itemize}[left=0pt]
  \item Normalize raw attendance rows into a JSON conference schema (conference metadata + entities).
  \item Build a matcher combining lexical keyword scoring and semantic similarity (embedding-backed) with structured filters (role, stage).
  \item Provide a web UI and programmatic API for interactive queries and explanations.
\end{itemize}

\subsection*{Repository and entry points}
\begin{itemize}[left=0pt]
  \item README: see README.md in repository root for install/run instructions.
  \item Importer CLI: \texttt{python3 -m conference\_matching.kaggle\_import} — normalizes raw CSV/ZIP into \texttt{data/conference\_kaggle.json}.
  \item Loader: \texttt{conference\_matching.data} exposes \texttt{load\_dataset()} and \texttt{normalize\_event\_attendance\_rows()}.
  \item Matcher: \texttt{conference\_matching.engine} builds the default hybrid matcher.
  \item Server: \texttt{server.py} starts the demo local server.
\end{itemize}

\subsection*{Design decisions}
Conservative normalization: attendees and events are constructed deterministically from attendance rows, deriving organization heuristically from email domain, and deriving topical sectors via phrase/token mapping. Hybrid ranking combines keyword overlap, concept matching, and optional semantic signals. Evaluation uses weak labels where each attendee's events form positives for retrieval experiments.

\section*{4 Data and Data Analysis}

\textbf{Data source:} Kaggle Event Attendance Dataset — https://www.kaggle.com/datasets/cankatsrc/event-attendance-dataset.

\textbf{Dataset Statistics:}
\begin{itemize}[left=0pt]
  \item Total Attendance Records: 200,000
  \item Unique Professional Attendees: 155,589
  \item Normalized Conference Sessions: 200
  \item Primary Tracks Identified: Artificial Intelligence, Customer Discovery, Community, etc.
\end{itemize}

\section*{5 Results and Evaluation}

\textbf{Evaluation approach:} We utilized the evaluation harness in \texttt{conference\_matching.evaluation}, which constructs synthetic queries from imported attendee profiles. For each attendee, their known session registrations are treated as ground-truth relevant items. We compared our hybrid matcher (which integrates lexical overlap, semantic concept similarity, and structured role/stage filtering) against a standard keyword-based lexical baseline.

\textbf{Quantitative Results} (weak-label evaluation, 3 attendee queries):

\begin{center}
\begin{tabular}{|l|c|c|}
\hline
\textbf{Metric} & \textbf{Hybrid} & \textbf{Keyword Baseline} \\
\hline
Precision@5 & 0.20 & 0.20 \\
Recall@5    & 1.00 & 1.00 \\
nDCG@5      & \textbf{1.00} & 0.54 \\
MRR         & \textbf{1.00} & 0.39 \\
\hline
\end{tabular}
\end{center}

The hybrid system (lexical + sentence-transformer embeddings via \texttt{all-MiniLM-L6-v2} + FAISS) achieves perfect nDCG@5 and MRR, ranking the correct session first for every query. The keyword-only baseline achieves the same recall but significantly lower ranking quality (MRR 0.39 vs.\ 1.00), a \textbf{157\% relative improvement} in MRR for the hybrid approach. Both systems achieve full Recall@5, confirming all relevant sessions are surfaced within the top 5 results.

\section*{6 Conclusions}

This implementation demonstrates that sparse attendance records can be successfully normalized into a rich entity store suitable for high-performance retrieval. By blending lexical, semantic, and structured features, we achieved a robust matching system that outperforms simple keyword search. Future work will focus on integrating large language models for more nuanced "reasoning" behind matches and expanding the ontology to cover more niche technical domains.

\section*{7 Submission Guidelines}

Include the PDF file, GCP endpoint link, and GitHub URL in your submission. Provide team members and repository on top of the document.

\section*{Appendix: Reproduce locally}

\begin{enumerate}[left=0pt]
  \item Create venv and install deps:
  \begin{verbatim}
python3 -m venv .venv
.venv/bin/python -m pip install -r requirements.txt
  \end{verbatim}
  \item Import the Kaggle dataset (example):
  \begin{verbatim}
python3 -m conference_matching.kaggle_import --input /path/to/event-attendance-dataset.zip
  \end{verbatim}
  \item Start server:
  \begin{verbatim}
python3 server.py
  \end{verbatim}
  \item Run evaluation:
  \begin{verbatim}
python3 -m conference_matching.evaluation
  \end{verbatim}
\end{enumerate}

\bigskip
\textit{(Replace placeholders and add experimental numbers and plots before final submission.)}

\end{document}
